%! Displaying and Analyzing Categorical Data — Unit 3, Lesson 3.1 — Homework (Answer Key)
% The lesson's SCORED (1 to 5) individual practice, assigned at Close & Assign and done outside
% class. DELTAMATH + JUSTIFICATIONS: DeltaMath carries items 1, 3, 4 (plain \blank here, \ans in
% the key; they stay as a worked reference). Items 2, 5, 6 are turned in: their \gradedblank
% boxes are copied verbatim to the key.
% THIRD CONTEXT: notes = getting to school at Oakdale High, Guided Practice = a library's
% streaming survey, homework = pet adoptions at a shelter.
% Items 3 and 4 are the CONTRAST PAIR: the same count (30), two different wholes.
% Item 5 is the crux head-on, with its justification. Item 2 is the interpret item.
% Item 6 is the spiral (Lesson 2.4's common scale) and a second justification.
% THE DATA (Maple Street Animal Shelter, 300 adoptions):
%             Dog  Cat  Total
%   18-34      90   60    150   60% 40%
%   35-54      55   45    100   55% 45%
%   55 and up  20   30     50   40% 60%
%   Total     165  135    300   55% 45%
%   30/50 = 60%; 30/135 = 22.2%; 165/300 = 55%; 60/150 = 40%
% POINTS: 12 in all. q2, q5, q6 rubric at 2 each (6); the DeltaMath score is scaled to the other 6.
% PAGE LOCKSTEP: the key differs only in -key for -boxes, the header, and \blank -> \ans; every
% \gradedblank and work block is byte-identical.
\documentclass[12pt]{article}
\usepackage{saar-article}
\usepackage{saar-key}   % pulls in -boxes; do NOT also load -boxes
\usepackage{grading}
\gradingassignment{saar-u03-l01-hw}
\gradingcourse{stats_analysis_algebraic_reasoning}
\newcommand{\TallMath}[1]{$\displaystyle #1\rule[-1.4em]{0pt}{3.2em}$}
% Handwriting room inside every \begin{work} block. Moves the blank and the
% key together, so the two can never drift apart.
\setlength{\workrowsep}{6pt}

\begin{document}

\pageheader{Unit 3, Lesson 3.1}{Homework --- Answer Key\strut}
% NAMESTRIP: no \namedateperiod here — mirrors the blank. NO teachernote here —
% teacher prose lives in the lesson plan.

\vspace{0.08in}

% Authored identically in homework/ and homework_key/.
\begin{remindbox}
\small
\textbf{This is your graded homework.} It is scored 1 to 5 and due at the start of the first
class after two study halls. It is assigned today and done on your own, outside class --- flip
back to the \emph{Keep in Mind} box on the cover and to rows 3 and 4 of your notes.
\end{remindbox}

\vspace{0.05in}

\boxguard[12]
\begin{scenariobox}[Maple Street Animal Shelter]{cerulean}
\small
Last year, $300$ pets were adopted from the shelter, each a dog or a cat: $150$ by adopters aged
$18$--$34$, $100$ by adopters aged $35$--$54$, and $50$ by adopters $55$ and over. Adopters
$18$--$34$ took home $90$ dogs and $60$ cats. Adopters $35$--$54$ took home $45$ cats. Adopters
$55$ and over took home $30$ cats.
\end{scenariobox}

\vspace{0.05in}

\boxguard
\begin{notesbox}{Pets and Their People}
\small
\textit{Write each final answer in its box. Only what is inside a box is graded.} Show your work
outside the boxes.
\begin{enumerate}[leftmargin=1.6em, itemsep=6pt, topsep=3pt]

  \item Finish the two-way table from the summary above. Subtract to find what it leaves out.

        \begin{center}
        \renewcommand{\arraystretch}{1.5}
        \begin{tabularx}{0.82\linewidth}{@{} l Y Y Y @{}}
        \toprule
        \textbf{Adopter age} & \textbf{Dog} & \textbf{Cat} & \textbf{Total} \\
        \midrule
        $18$--$34$ & $90$ & $60$ & $150$ \\
        $35$--$54$ & \ans{55} & $45$ & $100$ \\
        $55$ and over & \ans{20} & $30$ & $50$ \\
        \textbf{Total} & \ans{165} & \ans{135} & $300$ \\
        \bottomrule
        \end{tabularx}
        \end{center}

  \item What percent of all $300$ adoptions were dogs? Give the percent and say what it means
        for this shelter.

        \begin{work}
          \dfrac{165}{300} &= 55\%
        \end{work}

        \gradedresponse{\gradedblank[id=q2, type=rubric, points=2, display={55\%. Of the 300 pets adopted from the shelter last year, 55\% were dogs. This is a marginal relative frequency: the whole is every adoption.}]{\linewidth}{2in}}

\end{enumerate}
\end{notesbox}

\boxguard[30]
\begin{notesbox}{Pets and Their People, continued}
\small
\begin{enumerate}[leftmargin=1.6em, itemsep=5pt, topsep=2pt, start=3]

  \item Of the adopters $55$ and over, what percent took home a cat? \\[5pt]
        Denominator: \ans{50} \qquad
        Percent: \ans{60\%}

  \item Of the cats adopted, what percent went to adopters $55$ and over? \\[5pt]
        Denominator: \ans{135} \qquad
        Percent: \ans{22.2\%}

  \item The director says, \emph{``Adopters $18$--$34$ took home $60$ cats and adopters $55$ and
        over took home only $30$, so younger adopters are more likely to choose a cat.''} Do you
        agree? Justify your answer with the percents that settle it.
        \gradedresponse{\gradedblank[id=q5, type=rubric, points=2, display={Disagree. 60/150 = 40\% of adopters 18--34 chose a cat, but 30/50 = 60\% of adopters 55 and over did. The director compared counts from groups of different sizes; adopters 55 and over are more likely to choose a cat.}]{\linewidth}{2in}}

  \item \emph{Back to Lesson 2.4:} we compare groups only on a common scale. Which display should
        the shelter use to compare how likely each age group is to choose a cat? Justify your
        choice.

        {\centering
        \begin{tikzpicture}[x=1cm, y=0.0205cm]
          % Display A: side-by-side counts
          \node[font=\scriptsize\bfseries] at (2.6,112) {Display A: number of adoptions};
          \foreach \v in {30,60,90} { \draw[linegray] (0,\v) -- (5.4,\v); }
          \foreach \v in {0,30,60,90} { \node[left, font=\scriptsize] at (0,\v) {\v}; }
          \draw[thick] (0,0) -- (0,98);
          \draw[thick] (0,0) -- (5.4,0);
          \foreach \c/\a/\b/\g in {0.3/90/60/18--34, 2.1/55/45/35--54, 3.9/20/30/55+} {
            \fill[cerulean, draw=charcoal] (\c,0) rectangle ({\c+0.6},\a);
            \fill[goldbg, draw=charcoal] ({\c+0.6},0) rectangle ({\c+1.2},\b);
            \node[above, font=\tiny\bfseries, inner sep=1pt] at ({\c+0.3},\a) {\a};
            \node[above, font=\tiny\bfseries, inner sep=1pt] at ({\c+0.9},\b) {\b};
            \node[below, font=\scriptsize] at ({\c+0.6},0) {\g};
          }
          % Display B: 100% segmented
          \begin{scope}[xshift=7.1cm]
          \node[font=\scriptsize\bfseries] at (2.6,112) {Display B: percent of each age group};
          \foreach \v in {0,25,50,75,100} { \node[left, font=\scriptsize] at (0,{\v*0.98}) {\v\%}; }
          \draw[thick] (0,0) -- (0,98);
          \draw[thick] (0,0) -- (5.4,0);
          \foreach \c/\a/\b/\g in {0.5/60/40/18--34, 2.3/55/45/35--54, 4.1/40/60/55+} {
            \fill[cerulean, draw=charcoal] (\c,0) rectangle ({\c+0.9},{\a*0.98});
            \fill[goldbg, draw=charcoal] (\c,{\a*0.98}) rectangle ({\c+0.9},98);
            \node[font=\tiny\bfseries, text=white] at ({\c+0.45},{\a*0.49}) {\a\%};
            \node[font=\tiny\bfseries] at ({\c+0.45},{\a*0.98+\b*0.49}) {\b\%};
            \node[below, font=\scriptsize] at ({\c+0.45},0) {\g};
          }
          \end{scope}
          \node[font=\scriptsize] at (6.25,-42) {Dark bars: dogs.\quad Light bars: cats.};
        \end{tikzpicture}\par}

        \gradedresponse{\gradedblank[id=q6, type=rubric, points=2, display={Display B. It draws every age group as its own 100\%, so groups of different sizes sit on one common scale: 40\% of adopters 18--34 chose cats and 60\% of adopters 55 and over did. Display A shows counts, 60 and 30, and those depend on how big each group is.}]{\linewidth}{2in}}

\end{enumerate}
\end{notesbox}

\boxguard[5]   % the short spiralbox needs about five baselines; a larger guard risks a third page
\begin{spiralbox}
\small
\textbf{Next: Lesson 3.2, Probability Models.} A relative frequency is the first step toward a
probability: the percent of the time something happened is our estimate of how likely it is.
\end{spiralbox}

\end{document}
